\chapter{Reproductibilité et contrôles complémentaires}
\label{annexe:validation}
% Les intertitres de l'annexe coiffent des blocs de commandes, pas des
% developpements : ils n'ont ni numero ni entree dans la table des matieres,
% qui n'annonce ici que A.1 a A.4.
%
% Mecanisme : la classe numerote jusqu'au niveau 5 (secnumdepth=5) et n'inscrit
% dans la table que les niveaux 1 et 2. En abaissant secnumdepth a 2 ici, un
% \subsubsection (niveau 3) depasse les deux seuils : il perd son numero et
% n'entre pas dans la table. Son format est identique a celui d'un
% \subsection dans cette classe, donc le rendu ne change pas.
% \subsection* ne peut pas servir : la classe saisit l'argument avant de
% chercher l'etoile, et ses macros starsection et nostarsection sont
% identiques, donc l'etoile n'est jamais transmise.
\setcounter{secnumdepth}{2}
% Sans numero, un intertitre en graisse normale ne se distingue pas du texte :
% on le passe en gras. \origsubsubsection est la forme \@startsection que la
% classe enveloppe ; la redefinir ici conserve toute sa logique d'espacement.
\makeatletter
\renewcommand\origsubsubsection{\@startsection{subsubsection}{3}{\z@}%
                                     {12bp\@plus .2ex \@minus .2ex}%
                                     {.01pt \@plus .2ex}%
                                     {\normalfont\normalsize\bfseries\baselineskip=\simpleinter}}
\makeatother

Cette annexe recense les commandes et artefacts qui soutiennent les résultats du mémoire.
Les commandes sont exécutées depuis la racine du dépôt reproductible.
Le code source, les scripts de reproduction et les artefacts versionnés sont disponibles dans
le dépôt GitHub du projet~: \url{https://github.com/Barryguinea/memoirev3}.

% =====================================================================
\section{Environnement et commandes de reproduction}
% =====================================================================
\begin{verbatim}
python -m venv .venv
source .venv/bin/activate
python -m pip install --upgrade pip
python -m pip install -r requirements.txt
\end{verbatim}

L'environnement effectivement utilisé peut aussi être appelé avec le chemin absolu de son
interpréteur Python. Les résultats ne doivent pas être mélangés avec des artefacts produits
par une autre configuration du code.

La reproduction complète suit six étapes, dans l'ordre~:
\begin{enumerate}
  \item préparer l'environnement et placer \texttt{data/brut.csv} dans le projet~;
  \item exécuter l'évaluation et l'ablation du module HYPO~;
  \item exécuter l'analyse OFAT, la calibration exploratoire et la courbe
  détection-charge de HYPO~;
  \item exécuter le stress test HYPO à aire d'enveloppe appariée, puis la comparaison exploratoire des fusions~;
  \item exécuter l'analyse IceTag-SLS, l'audit bibliographique et régénérer les figures~;
  \item régénérer puis vérifier le manifeste SHA-256.
\end{enumerate}
Les commandes correspondantes sont détaillées ci-dessous.

\subsubsection{Évaluation et ablation du module HYPO}
\begin{verbatim}
python scripts/run_hypo_module_validation.py \
  --output-dir data/validation/hypo_module
\end{verbatim}

Cette commande régénère les 44 événements, les contrôles du protocole, le résumé par profil,
l'ablation et les tests appariés regroupés par vache.

\subsubsection{Stress test HYPO à aire d'enveloppe appariée}
\begin{verbatim}
python -m scripts.run_hypo_stress_validation
\end{verbatim}

La commande applique cinq formes moins favorables et un contrôle pas-seuls à trois positions
par vache. Le protocole est lu depuis \path{validation_hypo/stress_protocol.json}; il interdit
le réglage du détecteur et impose, pour chaque famille modifiée, la même aire discrète sous
l'enveloppe numérique que le profil graduel modéré de 48 heures. Les intervalles supprimés du
scénario de données manquantes sont exclus du calcul de cette aire. Cet appariement contrôle
l'intensité numérique programmée, sans prétendre égaliser l'amplitude physiologique ou
comportementale effectivement réalisée.

\subsubsection{Typologie des modes de défaillance}
\begin{verbatim}
python scripts/compute_failure_modes.py
\end{verbatim}

Cette commande relit les sorties événementielles de HYPO et de la fusion hiérarchique, puis
écrit les fréquences du Tableau~\ref{tab:failure_modes} dans
\path{hypo_module/failure_modes.csv}. Elle ne relance pas les détecteurs.

\subsubsection{Analyse OFAT du module HYPO}
\begin{verbatim}
python -m validation_hypo.sensitivity \
  --output-dir data/validation/hypo_module/parameter_sensitivity
\end{verbatim}

La référence et les 20 variantes utilisent les mêmes onze vaches, les mêmes fenêtres et
la même graine. L'analyse ne sélectionne pas de nouvelle configuration.

\subsubsection{Sensibilité des comparateurs à leur contamination}
\begin{verbatim}
python scripts/compute_comparator_sensitivity.py
\end{verbatim}

Cette commande rejoue l'ablation complète à cinq valeurs de \texttt{contamination}, de 0,02 à
0,15, sur les mêmes onze vaches et les mêmes 44 événements, puis écrit
\path{hypo_module/comparator_contamination_sensitivity.csv}. Elle contrôle que l'échec de
localisation d'Isolation Forest et de LOF ne dépend pas du réglage retenu pour eux. HYPO et le
comparateur pédométrique n'emploient pas Isolation Forest et servent de témoin~: leurs lignes
doivent rester constantes sur toute la plage.

\subsubsection{Sensibilité à la tolérance d'appariement des départs}
\begin{verbatim}
python scripts/compute_match_tolerance_sensitivity.py
\end{verbatim}

Cette commande rejoue la campagne principale à trois tolérances d'appariement, 0, 1 et
2~heures, et écrit \path{hypo_module/match_tolerance_sensitivity.csv}. Elle documente l'effet
de ce choix sur le taux de nouveau départ et vérifie que l'IoU reste inchangé, puisqu'il se
calcule sans apparier de départs. La constante
\texttt{DEFAULT\_MATCH\_TOLERANCE\_HOURS} de \path{validation_hypo/campaign.py} porte la
valeur de référence.

\subsubsection{Traçabilité des données (provenance CowAlert)}
\begin{verbatim}
python scripts/cowalert_provenance.py
\end{verbatim}

Cette commande confronte \path{data/brut.csv} à l'extraction
\path{data/cowalert_export.csv} relevée depuis le système CowAlert, et écrit les corrélations
et la proportion d'intervalles appariés dans \path{hypo_module/cowalert_provenance.csv}. Le
fichier source est confidentiel, demeure local et n'est pas versionné.

\subsubsection{Intervalles de confiance et tailles d'effet}
\begin{verbatim}
python scripts/compute_bootstrap_ci.py
\end{verbatim}

Ce script recalcule, à graine fixée, les intervalles de confiance à 95\,\% des métriques
attribuables (bootstrap regroupé par vache, 10\,000~rééchantillonnages), les tailles d'effet
de l'ablation et la différence appariée de charge de fond entre HYPO et le comparateur
pédométrique, puis les écrit dans \path{hypo_module/bootstrap_ci.csv}.

\subsubsection{Calibration individuelle et courbe détection-charge}
\begin{verbatim}
python scripts/per_cow_calibration.py
python scripts/detection_background_curve.py
\end{verbatim}

La première commande estime les facteurs uniquement sur la période de référence et conserve le résultat
négatif de l'expérience. La seconde réexécute neuf seuils de score sur les mêmes fenêtres,
écrit le tableau complet et produit la Figure~\ref{fig:detection_background_curve}. Aucune
des deux commandes ne remplace la configuration de référence.

\subsubsection{Comparaison exploratoire des fusions}
\begin{verbatim}
python -m validation_hybrid.sensitivity \
  --output-dir data/validation/hybrid_refined_full
\end{verbatim}

Cette commande exécute les cinq règles de fusion et quatre configurations de sensibilité sur
les mêmes onze vaches et les mêmes douze scénarios.

\subsubsection{Évaluation du temps d'exécution}
\begin{verbatim}
python scripts/benchmark_performance.py \
  --input data/brut.csv \
  --output-root data/performance \
  --fractions 1.0 --repeats 5 --collect-output \
  --canonical-json data/validation/performance_full_corpus.json
\end{verbatim}

Le fichier canonique retient la médiane des cinq exécutions complètes et les temps par étape.

\subsubsection{Analyse IceTag-SLS}
\begin{verbatim}
python -m validation_hybrid.mcgill_sls_validation
\end{verbatim}

Le répertoire source \texttt{mcgill\_iot\_cattle} doit être disponible selon le chemin
documenté dans le module. Aucune observation postérieure au 12 mars 2019 n'est utilisée.

\subsubsection{Figures du manuscrit}
\begin{verbatim}
python scripts/plot_parameter_sensitivity.py \
  --summary data/validation/hypo_module/parameter_sensitivity/\
parameter_sensitivity_summary.csv \
  --pdf memoire/figures/hypo_parameter_sensitivity_effects.pdf \
  --png memoire/figures/hypo_parameter_sensitivity_effects.png
python scripts/plot_manuscript_results.py
python scripts/plot_cow_admissibility.py
python scripts/compute_event_f1.py
python scripts/compute_channel_ablation.py
python scripts/compute_correlation_time.py
python scripts/compute_holm_adjustment.py
python scripts/audit_bibliography.py
python scripts/audit_manuscript_numbers.py
python scripts/update_validation_manifest.py
shasum -a 256 -c data/validation/validation_artifacts.sha256
\end{verbatim}

Le script bibliographique contrôle les 73 références~: 58 DOI sont confrontés aux métadonnées
Crossref ou DataCite et les 15 références sans DOI à une URL d'autorité préalablement revue.
Le premier script représente l'OFAT de HYPO. Le second lit les CSV finaux et régénère les
figures de fusion, de scénarios et de SLS; il reconstruit également quatre trajectoires
représentatives pour l'inspection visuelle. Les quatre scripts suivants écrivent les métriques
F1, l'ablation par famille, les temps d'autocorrélation et les valeurs de $p$ corrigées par
Holm dans \path{data/validation/derived_metrics}. Le script d'audit compare ensuite 398 valeurs
empiriques du manuscrit à leurs CSV ou JSON sources. Les deux dernières commandes reconstruisent
le manifeste des artefacts scientifiques puis en vérifient chaque empreinte.

\subsubsection{Tests et interface}
\begin{verbatim}
pytest -q
python -m streamlit run app.py
\end{verbatim}

L'interface est ensuite accessible à \texttt{http://localhost:8501}. Les tests vérifient la
cohérence logicielle, non la validité clinique.

\subsubsection{Compilation}
\begin{verbatim}
cd memoire
latexmk -C
latexmk -lualatex -interaction=nonstopmode -halt-on-error main.tex
\end{verbatim}

Le nettoyage préalable n'est pas facultatif. La conformité PDF/A-1b est obtenue avec
\texttt{pdfx}, qui charge lui-même \texttt{hyperref}~; sur une recompilation effectuée sans
nettoyage, la première passe relit les fichiers auxiliaires de la compilation précédente avant
que cette machinerie ne soit disponible et s'interrompt sur une séquence de contrôle non
définie. La compilation à froid, elle, aboutit sans avertissement.

% =====================================================================
\section{Règles d'interprétation et unité statistique}
% =====================================================================
\subsubsection{Règles d'interprétation}
Les formulations suivantes sont exclues des conclusions:
\begin{itemize}
  \item appeler « sensibilité clinique » le taux de détection des injections;
  \item appeler « faux positifs » les notifications produites sur des données sans labels;
  \item interpréter une instabilité isolée comme une boiterie précoce confirmée;
  \item attribuer l'association SLS au système sans mentionner le groupe \textit{Exercise};
  \item présenter les paramètres techniques comme des seuils physiologiques optimaux;
  \item remplacer le seuil de référence par 0,14 à partir de la courbe descriptive, sans
  jeu de développement indépendant;
  \item comparer directement 87,9\,\% de couverture HYPO et 72,7\,\% de détection
  nécessitant une vérification dans l'extension hybride, car les critères et les ensembles de scénarios diffèrent;
  \item attribuer l'écart entre 0,402 et 0,446 au seul algorithme, car le délai réfractaire
  passe de 24 à 12 heures dans la comparaison des fusions;
  \item traiter les intervalles de quinze minutes comme des sujets indépendants.
\end{itemize}

Ces règles font partie du protocole de reproductibilité: reproduire les calculs sans en
respecter la portée conduirait à une conclusion scientifiquement différente.

\subsubsection{Unité statistique}
Les 132 lignes événementielles ne représentent pas 132 vaches indépendantes. Elles sont
regroupées dans onze vaches. Les résultats par scénario sont descriptifs; toute comparaison
inférentielle doit agréger préalablement par vache. Dans l'analyse SLS, chaque vache fournit
une seule valeur par mesure sur la fenêtre de sept jours.

La campagne HYPO comporte 44 événements mais repose sur les mêmes onze vaches. Les tests de
l'ablation utilisent donc les moyennes par vache, jamais les 44 lignes comme observations
indépendantes.

Le stress test comporte 198 événements physiques et 990 lignes après application des cinq
variantes. Les trois positions et les six formes restent imbriquées dans les mêmes onze
vaches. Les intervalles de confiance rééchantillonnent donc les vaches, et les tests exacts
appariés portent sur onze différences moyennes par vache.

% =====================================================================
\section{Données et artefacts}
% =====================================================================
\subsubsection{Artefacts principaux}
Le Tableau~\ref{tab:artefacts_validation} liste les artefacts qui soutiennent les résultats.

\begingroup
\small
\begin{longtable}{p{0.43\textwidth}p{0.47\textwidth}}
  \caption{Artefacts qui soutiennent les résultats.}
  \label{tab:artefacts_validation}\\
  \toprule
  Fichier & Contenu \\
  \midrule
  \endfirsthead
  \toprule Fichier & Contenu \\ \midrule \endhead
  \path{hypo_module/events_primary.csv} & 44 événements de la campagne principale HYPO. \\
  \path{hypo_module/protocol_checks.csv} & Contrôles automatiques de la campagne HYPO. \\
  \path{hypo_module/detection_summary.csv} & Résultats attribuables par profil. \\
  \path{hypo_module/ablation_primary.csv} & Sorties des cinq variantes sur les mêmes fenêtres. \\
  \path{hypo_module/ablation_summary.csv} & Résumé de l'ablation HYPO, IF, LOF et pédométrique. \\
  \path{hypo_module/ablation_tests_by_cow.csv} & Tests appariés après agrégation par vache. \\
  \path{hypo_module/parameter_sensitivity/parameter_sensitivity_summary.csv} & OFAT des dix paramètres HYPO. \\
  \path{hypo_module/bootstrap_ci.csv} & IC95 (bootstrap par vache), tailles d'effet et comparaison appariée de la charge de fond. \\
  \path{hypo_module/failure_modes.csv} & Fréquences reproductibles des modes de défaillance. \\
  \path{hypo_module/per_cow_calibration.csv} & Comparaison agrégée avant et après calibration. \\
  \path{hypo_module/per_cow_calibration_by_cow.csv} & Facteurs et charges de fond par vache. \\
  \path{hypo_module/cowalert_provenance.csv} & Concordance IceQube-CowAlert (corrélations, identité). \\
  \path{hypo_module/detection_background_curve/detection_background_curve.csv} & Neuf points de la courbe détection-charge. \\
  \path{hypo_stress/events.csv} & Sorties des cinq variantes sur les événements de stress à aire d'enveloppe appariée. \\
  \path{hypo_stress/scenario_summary.csv} & Résultats HYPO par forme moins favorable. \\
  \path{hypo_stress/paired_comparisons.csv} & Tests exacts appariés après agrégation par vache. \\
  \path{hypo_stress/summary.json} & Protocole, empreinte et contrôle d'intégrité des aires d'enveloppe. \\
  \path{derived_metrics/event_f1.csv} & Précision, rappel, F1 et IC95 selon trois définitions événementielles. \\
  \path{derived_metrics/channel_ablation_summary.csv} & Synthèse de l'ablation par famille de signaux. \\
  \path{derived_metrics/channel_ablation_events.csv} & Sorties événementielles détaillées de l'ablation par famille. \\
  \path{derived_metrics/correlation_time.csv} & Temps d'autocorrélation et tailles effectives aux deux grains d'analyse. \\
  \path{input_provenance.json} & Empreinte, taille et dimensions du corpus brut confidentiel. \\
  \path{manuscript_number_audit.csv} & Correspondance entre 398 valeurs du manuscrit et leurs artefacts sources. \\
  \path{performance_full_corpus.json} & Médiane et décomposition du temps d'exécution de la chaîne complète. \\
  \path{hybrid_refined_full/comparison_summary.csv} & Comparaison des fusions et variantes de sensibilité. \\
  \path{hybrid_refined_full/events_fusion_hierarchical.csv} & 132 évaluations de la fusion hiérarchique exploratoire. \\
  \path{hybrid_refined_full/events_fusion_hypo_only.csv} & Résultats de la branche HYPO seule. \\
  \path{hybrid_refined_full/events_fusion_instability_only.csv} & Résultats de la branche INSTABILITÉ seule. \\
  \path{hybrid_refined_full/events_fusion_or.csv} & Résultats de la fusion OU. \\
  \path{hybrid_refined_full/events_fusion_sequential_24_72h.csv} & Résultats de la fusion séquentielle. \\
  \path{hybrid_refined_full/events_parameter_sensitivity_*.csv} & Évaluations des variantes de paramètres. \\
  \path{mcgill_sls/mcgill_cohort_all_variants.csv} & Cohorte SLS et sorties des cinq fusions. \\
  \path{mcgill_sls/mcgill_exclusions.csv} & Motifs d'exclusion et couverture. \\
  \path{mcgill_sls/mcgill_metrics.csv} & AUC, Mann-Whitney et Spearman. \\
  \path{mcgill_sls/mcgill_treatment_sensitivity.csv} & Tests exacts globaux et stratifiés par traitement. \\
  \path{mcgill_sls/mcgill_multiplicity_sensitivity.csv} & Sensibilité max-stat à la sélection parmi cinq ou vingt analyses SLS. \\
  \path{mcgill_sls/mcgill_summary.json} & Protocole et synthèse de l'analyse SLS. \\
  \path{bibliography_audit.csv} & Audit des 73 références par DOI ou URL d'autorité. \\
  \path{memoire/figures/manual_review.png} & Planche reproductible de l'inspection visuelle ciblée. \\
  \bottomrule
\end{longtable}
\endgroup

Sauf pour la planche d'inspection, les chemins du tableau sont relatifs à
\texttt{data/validation/}.

\subsubsection{Intégrité des fichiers}
Le fichier \texttt{data/validation/validation\_artifacts.sha256} contient les empreintes SHA-256
des sorties qui alimentent ce mémoire~: performance attribuable, ablation, analyses
exploratoires, temps d'exécution, comparaison des fusions et concordance IceTag-SLS. Il scelle
aussi \texttt{input\_provenance.json}, qui consigne l'empreinte du corpus brut confidentiel sans
en révéler le contenu. Les chemins y sont relatifs à la racine du projet. Après une reproduction,
la commande suivante vérifie le manifeste~:

\begin{verbatim}
shasum -a 256 -c data/validation/validation_artifacts.sha256
\end{verbatim}

Toute modification des scripts, des données, des paramètres ou de l'environnement exige une
nouvelle exécution et un nouveau manifeste. Le PDF seul ne constitue pas une trace suffisante
de l'expérience.

% =====================================================================
\section{Profils d'injection et variabilité}
% =====================================================================
\subsubsection{Amplitudes des profils d'injection}
Les Tableaux~\ref{tab:amplitudes_hypo} et~\ref{tab:amplitudes_extension} donnent les
changements relatifs appliqués aux signaux bruts. Les amplitudes des trois baisses sont
communes aux deux campagnes, mais leurs durées diffèrent~: 36 à 60 heures dans la campagne
principale HYPO, contre 24 à 48 heures dans l'extension bidirectionnelle. Les modifications
suivent une enveloppe progressive. Seule l'amplitude modérée des pas ($-29$\,\%) est ancrée
sur un ordre de grandeur publié \citep{paudyal2021lying}~; les autres valeurs sont des
hypothèses de test et non des estimations cliniques.

\begin{table}[H]
  \centering
  \caption{Profils de la campagne principale HYPO.}
  \label{tab:amplitudes_hypo}
  \small
  \begin{tabular}{lrrrrr}
    \toprule
    Profil & Durée & Pas & Mouvement & Transitions & Posture \\
    \midrule
    Graduel léger & 36 h & $-12$\,\% & $-10$\,\% & $-12$\,\% & $+3$\,\% \\
    Graduel modéré & 48 h & $-29$\,\% & $-25$\,\% & $-25$\,\% & $+7$\,\% \\
    Graduel marqué & 60 h & $-42$\,\% & $-38$\,\% & $-40$\,\% & $+12$\,\% \\
    Variation brève & 3 h & $+15$\,\% & $\cdot$ & $\cdot$ & $\cdot$ \\
    \bottomrule
  \end{tabular}
\end{table}

\begin{table}[H]
  \centering
  \caption{Profils bidirectionnels et contrôles de confusion.}
  \label{tab:amplitudes_extension}
  \small
  \begin{tabular}{lrrrrr}
    \toprule
    Profil & Durée & Pas & Mouvement & Transitions & Posture/osc. \\
    \midrule
    \multicolumn{6}{l}{\textit{Hypoactivité de l'extension}} \\
    Léger & 24 h & $-12$\,\% & $-10$\,\% & $-12$\,\% & $+3$/$\cdot$\,\% \\
    Modéré & 36 h & $-29$\,\% & $-25$\,\% & $-25$\,\% & $+7$/$\cdot$\,\% \\
    Marqué & 48 h & $-42$\,\% & $-38$\,\% & $-40$\,\% & $+12$/$\cdot$\,\% \\
    \addlinespace
    \multicolumn{6}{l}{\textit{Instabilité (hausses disproportionnées aux pas)}} \\
    Léger & 8 h & $0$\,\% & $+25$\,\% & $+30$\,\% & $+6/+8$\,\% \\
    Modéré & 12 h & $+5$\,\% & $+45$\,\% & $+50$\,\% & $+10/+12$\,\% \\
    Marqué & 16 h & $+8$\,\% & $+65$\,\% & $+70$\,\% & $+15/+18$\,\% \\
    \addlinespace
    \multicolumn{6}{l}{\textit{Contrôles et confondants}} \\
    Pic capteur & 15 min & $\cdot$ & $+200$\,\% & $\cdot$ & $\cdot$ \\
    Exercice bref & 2 h & $+45$\,\% & $+45$\,\% & $+45$\,\% & $\cdot$/$+5$\,\% \\
    Activité de type œstrus & 6 h & $+40$\,\% & $+40$\,\% & $+40$\,\% & $\cdot$/$+10$\,\% \\
    Hypoactivité non locomotrice & 24 h & $-30$\,\% & $-28$\,\% & $-25$\,\% & $+8$/$\cdot$\,\% \\
    Manipulation & 1 h & $+10$\,\% & $+55$\,\% & $+70$\,\% & $+18/+20$\,\% \\
    \bottomrule
  \end{tabular}
\end{table}

Dans la dernière colonne du Tableau~\ref{tab:amplitudes_extension}, la première valeur est
le transfert ou l'oscillation posturale et la seconde l'oscillation générale appliquée aux
signaux. Un point ($\cdot$) indique l'absence de modification correspondante.

Le profil séquentiel enchaîne une instabilité ($+50$\,\% de mouvement, $+55$\,\% de
transitions, $+12/+14$\,\% de posture/oscillation sur 12~h), un délai de 24~h, puis une
hypoactivité ($-32$\,\% de pas, $-28$\,\% de mouvement, $-30$\,\% de transitions et
$+8$\,\% de transfert postural sur 36~h).

\subsubsection{Invariants des injections}
\begin{enumerate}
  \item chaque scénario est injecté dans une copie distincte des données;
  \item le début est strictement postérieur à la période de référence;
  \item les pas, le mouvement et les transitions restent non négatifs;
  \item le temps couché plus le temps debout est conservé;
  \item les enveloppes sont progressives et ne contiennent pas d'alternance extrême;
  \item chaque injection modifie effectivement au moins un signal attendu, la fenêtre étant
  glissée si nécessaire pour éviter une perturbation sans effet sur un signal nul;
  \item l'exécution propre de la même vache sert de référence d'attribution;
  \item le fichier \texttt{data/brut.csv} n'est jamais modifié.
\end{enumerate}

\subsubsection{Variabilité inter-vache du fond}
Le taux de fond agrégé (0,402 notification par vache-jour pour HYPO) masque une variabilité
individuelle importante. Sur les onze vaches admissibles (45 notifications, 112~vaches-jours),
le taux varie de 0,197 à 0,590 selon la vache, soit un facteur d'environ trois
(Tableau~\ref{tab:fond_par_vache}).

\begin{table}[H]
  \centering
  \caption{Distribution du fond de notifications HYPO par vache (campagne principale).}
  \label{tab:fond_par_vache}
  \small
  \begin{tabular}{lccc}
    \toprule
    Groupe & $n$ vaches & Notifications ($\approx$10~j) & Taux/vache-jour \\
    \midrule
    Stable ($\leq 0{,}30$) & 2 & 2--3 & 0,197--0,295 \\
    Intermédiaire (0,30--0,45) & 6 & 4 & 0,389--0,393 \\
    Active ($> 0{,}45$) & 3 & 5--6 & 0,491--0,590 \\
    \midrule
    Ensemble & 11 & 45 (total) & 0,402 (agrégé) \\
    \bottomrule
  \end{tabular}
\end{table}

Cette dispersion décrit une hétérogénéité de charge, mais elle ne démontre pas que le niveau
d'activité moyen en est la cause. Une calibration simple fondée sur la dispersion du score de
période de référence a été testée sans réduire le fond ni son coefficient de variation, et elle a dégradé
IoU20. Elle n'est donc pas retenue. Une méthode future devrait être définie sur un jeu de
développement distinct avant d'être évaluée. Ces chiffres décrivent une charge opérationnelle
et ne constituent pas un taux de faux positifs, faute de labels cliniques.
